\documentclass[a4paper]{article}
\usepackage{amsfonts}
\usepackage{amsmath}
\usepackage{amssymb}
\usepackage{graphicx}     

\begin{document}
  
\noindent \large Auteur: Rubin Cuypers \\
\noindent \large Studierichting: Informatica\\
\noindent \large Oefeningenreeks 3C nr. 14.2\\

\medskip

\normalsize
zij $h>0$\\

$ f:[a-b,a+b] \rightarrow \mathbb{R}$ continu en differtieerbaar op $ ]a-b,a+b[$\\

$\exists \theta_2 \in ]0,1[ : \dfrac{f(a+h)-2f(a)+f(a-h)}{h}=f'(a+\theta_2h)-f'(a-\theta_2h)$\\

$F:[0,1] \rightarrow \mathbb{R} : x \rightarrow f(a+hx)+f(a-hx)$\\

$F$ is continu op $[0,1]$ en differentieerbaar op $ ]0,1[ $ want het is een samenstelling van $f$ met 2 lineaire functies $a+hx$ en $a-hx$.\\

$\Rightarrow$ Lagrange\\

(i) $ F:[0,1] \rightarrow \mathbb{R} $ is continu op $[0,1]$\\

(ii) en differentieerbaar op $]0,1[$\\

$\Rightarrow \exists \theta_2 \in ]0,1[ : F'(\theta_2) =\dfrac{F(1)-F(0)}{1}$\\

$F'(\theta_2)=h(f'(a+h\theta_2)-f'(a-h\theta_2))$\\

$\Rightarrow h(f'(a+h\theta_2)-f'(a-h\theta_2))=f(a+h)+f(a-h)-(f(a)+f(a))$\\

$\Rightarrow f'(a+h\theta_2)-f'(a-h\theta_2)=\dfrac{f(a+h)+f(a-h)-2f(a)}{h}$\\


\begin{thebibliography}{999}
%\bibitem{a} Referentie
\end{thebibliography}
\end{document} 
